%% ──────────────────────────────────────────────
%%  پیوست پ – ایران پس از ۱۳۵۷
%%  فایل: appendices/app-c-iran-post1979.tex
%% ──────────────────────────────────────────────
\chapter{ایران پس از ۱۳۵۷: اسناد، عملیات‌ها، و ابزارهای مداخله}
\label{app:iran-post1979}

\begin{tcolorbox}[mybox, title={درباره‌ی این پیوست}]
این پیوست به بررسی اَشکال مختلف مداخله‌ی خارجی در ایران پس از انقلاب اسلامی ۱۳۵۷ می‌پردازد. تمرکز بر اسناد منتشرشده، عملیات‌های مستند، و ابزارهای شناخته‌شده است. از ارائه‌ی ادعاهای بدون مستند پرهیز شده و تنها به منابع معتبر و قابل‌دسترس عمومی استناد می‌شود.
\end{tcolorbox}

%% ================================================================
\section{اسناد منتشرشده‌ی سیا درباره‌ی ایران}
\label{sec:app-c-cia}
%% ================================================================

\needspace{6\baselineskip}
آرشیو آزادشده‌ی سیا (\lat{CIA FOIA Reading Room}) شامل هزاران سند درباره‌ی ایران است. مهم‌ترین مجموعه‌ها:

\begin{enumerate}
\item \textbf{اسناد کودتای ۲۸ مرداد ۱۹۵۳:} سیا در سال ۲۰۱۷ اسناد کاملی از عملیات \lat{TPAJAX} منتشر کرد که جزئیات برنامه‌ریزی، اجرا، و ارزیابی کودتا علیه مصدق را نشان می‌دهد.\footnote{\lr{Central Intelligence Agency. \textit{Declassified Documents Reference System: Iran 1953 Coup}, released 2017, available at \url{https://www.cia.gov/readingroom/}.}}

\item \textbf{ارزیابی‌های اطلاعاتی ملی (\lat{NIE}) درباره‌ی ایران:} ده‌ها ارزیابی از دهه‌ی ۱۹۸۰ تاکنون، شامل تحلیل توانمندی‌های نظامی، برنامه‌ی هسته‌ای، و ساختار قدرت.\footnote{\lr{National Intelligence Estimates on Iran, various years. Partially declassified. Available through the National Security Archive, George Washington University.}}

\item \textbf{اسناد دوران گروگان‌گیری (۱۹۷۹–۱۹۸۱):} شامل ارزیابی‌های روزانه و طرح‌های نجات (عملیات \concept{پنجه‌ی عقاب} / \lat{Eagle Claw}).\footnote{\lr{Bowden, Mark. \textit{Guests of the Ayatollah: The Iran Hostage Crisis}. New York: Grove Press, 2006.}}

\item \textbf{اسناد ایران-کنترا (۱۹۸۵–۱۹۸۷):} فروش مخفیانه‌ی سلاح به ایران (علی‌رغم تحریم‌ها) و استفاده از درآمدش برای تأمین مالی شورشیان نیکاراگوا.\footnote{\lr{Walsh, Lawrence E. \textit{Iran-Contra: The Final Report}. New York: Times Books, 1994.}}
\end{enumerate}

%% ================================================================
\section{عملیات‌های مستند علیه جمهوری اسلامی}
\label{sec:app-c-operations}
%% ================================================================

\needspace{6\baselineskip}

\begin{table}[htbp]
\centering
\caption{عملیات‌های مستند علیه جمهوری اسلامی ایران}
\label{tab:app-c-operations}
\begin{adjustbox}{max width=\linewidth}
\begin{tabularx}{\linewidth}{
    >{\centering\arraybackslash}p{1.5cm}
    >{\raggedleft\arraybackslash}p{46mm}
    >{\raggedleft\arraybackslash}p{35mm}
    >{\raggedleft\arraybackslash}X
}
\toprule
\rowcolor{PrimaryDeep}
\textcolor{white}{\textbf{سال}} &
\textcolor{white}{\textbf{عملیات/رویداد}} &
\textcolor{white}{\textbf{عامل}} &
\textcolor{white}{\textbf{منبع}} \\
\midrule
۱۹۸۰ & حمایت از کودتای نوژه & آمریکا/ارتش شاهنشاهی & اسناد آزادشده‌ی سیا \\
\rowcolor{NeutralLight}
۱۹۸۰–۸۸ & حمایت از عراق در جنگ & آمریکا/غرب/عرب & مستندات گسترده \\
۱۹۸۸ & سرنگونی پرواز ۶۵۵ ایران‌ایر & ناو آمریکایی وینسنز & مستند — ۲۹۰ کشته \\
\rowcolor{NeutralLight}
۲۰۱۰ & حمله‌ی سایبری استاکس‌نت & آمریکا/اسرائیل & تأیید ضمنی اوباما \\
۲۰۱۰–۲۰ & ترور دانشمندان هسته‌ای & اسرائیل (منتسب) & گزارش‌های متعدد رسانه‌ای \\
\rowcolor{NeutralLight}
۲۰۱۸ & خروج از برجام + فشار حداکثری & آمریکا (ترامپ) & اسناد رسمی \\
۲۰۲۰ & ترور ژنرال سلیمانی & آمریکا & تأیید رسمی پنتاگون \\
\bottomrule
\end{tabularx}
\end{adjustbox}
\end{table}

%% ================================================================
\section{تحریم‌ها به‌مثابه‌ی ابزار مداخله}
\label{sec:app-c-sanctions}
%% ================================================================

\needspace{6\baselineskip}
تحریم‌های ایران از ۱۹۷۹ آغاز شده و طی چهار دهه لایه‌لایه گسترش یافته‌اند:

\begin{enumerate}
\item \textbf{تحریم‌های اولیه (۱۹۷۹–۱۹۹۵):} مسدودسازی دارایی‌ها، ممنوعیت تجارت مستقیم
\item \textbf{قانون داماتو (۱۹۹۶):} تحریم‌های ثانویه — شرکت‌های غیرآمریکایی نیز هدف
\item \textbf{تحریم‌های شورای امنیت (۲۰۰۶–۲۰۱۰):} قطعنامه‌های ۱۷۳۷، ۱۷۴۷، ۱۸۰۳، ۱۹۲۹
\item \textbf{تحریم‌های فلج‌کننده (۲۰۱۲):} تحریم بانک مرکزی و صادرات نفت — تورم بالای ۴۰٪
\item \textbf{برجام و تعلیق نسبی (۲۰۱۶–۲۰۱۸):} رفع بخشی از تحریم‌ها
\item \textbf{فشار حداکثری ترامپ (۲۰۱۸–۲۰۲۱):} بازگشت همه‌ی تحریم‌ها + تحریم‌های جدید
\end{enumerate}

\begin{tcolorbox}[enemybox, title={اثر تحریم‌ها بر مردم ایران}]
\begin{itemize}
\item کاهش ارزش ریال: از ۱۲٬۰۰۰ تومان (۲۰۱۸) به بیش از ۶۰٬۰۰۰ تومان (۲۰۲۴) برای هر دلار
\item تورم مزمن: میانگین بالای ۳۰٪ سالانه
\item کمبود دارو: علی‌رغم «استثنای بشردوستانه»، عملاً بانک‌ها از تراکنش‌های دارویی خودداری می‌کنند
\item فرار مغزها: مهاجرت گسترده‌ی نخبگان
\item فقر: افزایش قابل‌توجه خط فقر
\item \textbf{اما:} تحریم‌ها رژیم را تغییر نداده‌اند — بلکه ناامیدی آموخته‌شده‌ی مردم را تشدید کرده‌اند
\end{itemize}
\end{tcolorbox}

%% ================================================================
\section{جنگ سایبری}
\label{sec:app-c-cyber}
%% ================================================================

\needspace{6\baselineskip}
\begin{itemize}
\item \textbf{استاکس‌نت (۲۰۱۰):} نخستین «سلاح سایبری» جهان — ویروسی که سانتریفیوژهای نطنز را فیزیکاً تخریب کرد. عملیات مشترک آمریکا و اسرائیل (نام رمز: \concept{بازی‌های المپیک} / \lat{Olympic Games}).\footnote{\lr{Sanger, David E. \textit{Confront and Conceal: Obama's Secret Wars and Surprising Use of American Power}. New York: Crown, 2012, pp. 188–225.}}
\item \textbf{Duqu (۲۰۱۱):} بدافزار جاسوسی مرتبط با استاکس‌نت
\item \textbf{Flame (۲۰۱۲):} بدافزار جاسوسی پیچیده — جمع‌آوری اطلاعات از رایانه‌های ایرانی
\item \textbf{حملات مداوم:} ده‌ها حمله‌ی سایبری گزارش‌شده به زیرساخت‌های ایران از ۲۰۱۰ تاکنون
\end{itemize}

%% ================================================================
\section{فعالیت رسانه‌های فارسی‌زبان خارجی}
\label{sec:app-c-media}
%% ================================================================

\needspace{6\baselineskip}

\begin{table}[htbp]
\centering
\caption{رسانه‌های فارسی‌زبان خارجی و منابع مالی‌شان}
\label{tab:app-c-media}
\begin{adjustbox}{max width=\linewidth}
\begin{tabularx}{\linewidth}{
    >{\raggedleft\arraybackslash}p{3cm}
    >{\raggedleft\arraybackslash}p{5cm}
    >{\centering\arraybackslash}p{1.5cm}
    >{\raggedleft\arraybackslash}X
}
\toprule
\rowcolor{PrimaryDeep}
\textcolor{white}{\textbf{رسانه}} &
\textcolor{white}{\textbf{تأمین مالی}} &
\textcolor{white}{\textbf{تأسیس}} &
\textcolor{white}{\textbf{ماهیت}} \\
\midrule
بی‌بی‌سی فارسی & دولت بریتانیا & ۱۹۴۰ & خبری-تحلیلی \\
\rowcolor{NeutralLight}
صدای آمریکا فارسی & دولت آمریکا (USAGM) & ۱۹۴۲ & خبری-تحلیلی \\
رادیو فردا & دولت آمریکا (RFE/RL) & ۲۰۰۲ & خبری-تحلیلی \\
\rowcolor{NeutralLight}
ایران اینترنشنال & سعودی (ولنتس) & ۲۰۱۷ & خبری-سیاسی \\
من‌وتو & خصوصی (لندن) & ۲۰۱۰ & سرگرمی-خبری \\
\bottomrule
\end{tabularx}
\end{adjustbox}
\end{table}

بدون قضاوت ارزشی درباره‌ی محتوای این رسانه‌ها، نکته‌ی ساختاری این است: \textbf{اکثر رسانه‌های فارسی‌زبان پرمخاطب، تأمین مالی خارجی (دولتی یا شبه‌دولتی) دارند.} این واقعیت — مستقل از کیفیت محتوا — آن‌ها را در چارچوب «تسهیل‌گری فرهنگی مداخله» (فصل~\ref{ch:missionary}) قرار می‌دهد.

%% ================================================================
\section{جمع‌بندی}
\label{sec:app-c-conclusion}
%% ================================================================

ایران پس از ۱۳۵۷ هدف تقریباً تمامی اَشکال مداخله‌ی خارجی — \textit{به‌جز} مداخله‌ی نظامی مستقیم — بوده است: کودتای نافرجام، حمایت از دشمن در جنگ، تحریم‌های فلج‌کننده، ترور دانشمندان، جنگ سایبری، و فعالیت رسانه‌ای گسترده. دلیل نبود مداخله‌ی نظامی مستقیم ترکیبی از عوامل بازدارنده است: جمعیت بزرگ، جغرافیای دشوار، توانمندی نظامی نسبی (موشکی و نیروی دریایی)، حمایت روسیه و چین، و تجربه‌ی شکست مداخلات مشابه در منطقه (عراق، افغانستان).

اما نکته‌ی اصلی این است: \textbf{عدم مداخله‌ی نظامی مستقیم به معنای عدم مداخله نیست.} مداخله در ایران شکل «نرم» اما مداوم و ساختاری دارد — و اثرات آن بر اقتصاد، جامعه، و روان‌شناسی جمعی ایرانیان عمیق و بلندمدت است. همان‌گونه که در فصل~\ref{ch:social-psychology} بحث شد، تحریم‌های طولانی‌مدت ناامیدی آموخته‌شده را تشدید و عاملیت جمعی را تضعیف می‌کنند — و بدین ترتیب پارادوکسی ایجاد می‌شود: ابزاری که قرار است «رژیم را تغییر دهد»، در عمل \textit{مردم} را تضعیف و رژیم را — که به ابزارهای سرکوب دسترسی دارد — نسبتاً دست‌نخورده باقی می‌گذارد.

\cleardoublepage